\documentclass[12 pt]{article}
\usepackage[a4paper, top=2cm, bottom=2cm, left=2cm, right=2cm]{geometry}
\usepackage{amsmath}
\usepackage{siunitx}
\usepackage{physics}
\usepackage{gensymb}
\usepackage{float}         
\usepackage{booktabs}
\usepackage{tabularx}
\usepackage{graphicx} 
\usepackage{caption}  

\usepackage{listings}
\usepackage{xcolor}

\date{08/2026}
\title{\textbf{\LARGE Four tap FIR Low Pass Filter} \\ RTL-GDSII  \vspace{8cm}}
\author{ Julia E. Benítez }  

\begin{document}
	\newgeometry{top=10.5cm, bottom=2cm, left=2cm, right=2cm}
	\maketitle
	% ---
	\newpage
	\newgeometry{top=2cm, bottom=2cm, left=2cm, right=2cm}
	
	\section{Objective}
	Design, simulate, and analyze the behavior of a 4-tap FIR Low Pass Filter as an application example of the complete RTL-to-GDSII ASIC design flow using Synopsys EDA tools. \\ 
	
\textit{	NOTE: The filter is intentionally limited to four taps in order to reduce hardware complexity while still demonstrating the complete digital implementation.}
	
	\section{Specifications}
	
	\begin{table}[H]
		
		\centering
		\label{tab:fir_specs}
		\begin{tabular}{ll}
			\toprule
			\textbf{Parameter} & \textbf{Specification} \\
			\midrule
			Filter type & Low-pass FIR \\
			Number of taps & 4 \\
			Design method & Hamming window \\
			Sampling frequency ($f_s$) & $100\,\mathrm{kHz}$ \\
			Nyquist frequency ($f_N$)& $50\,\mathrm{kHz}$ \\
			Cutoff frequency ($f_c$) & $12.5\,\mathrm{kHz}$ \\
			Normalized cutoff frequency & $0.25\,f_N$ ($0.125\,f_s$) \\
			Input word length & 16-bit signed fixed-point \\
			Coefficient word length & 16-bit signed fixed-point \\
			Technology & GlobalFoundries 180\,nm Standard Cell Library \\
			Implementation flow & RTL $\rightarrow$ GDSII ASIC Flow \\
			\bottomrule
		\end{tabular}
			\caption{Design specifications of the proposed FIR filter.}
	\end{table}
	
	\section{Design and implementation}
	
	\subsubsection{Theoretical filter design}
	A Filter is a simple device that allows altering the frequency content of a signal, and can be represented as mathematical expression with certain properties. For digital signals the equivalent model is established by using discrete time functions. 
	
	Properties of filters are better understood as functions of frequency domain. This functions doesn't introduce new frequencies meaning that they are linear time-invariant systems.
	
	 By function, they are classified as high pass, low pass, band stop, or band pass. In this case, the focus remains in the low pass filter which purpose is keep frequencies below certain value, called cutoff frequency , and attenuated or suppress the values above it.   
	
	To obtain a good model function for such task and still be a realizable design, the function must be causal or, mathematically, satisfy the Paley–Wiener Theorem. This implies that the ideal filter is not realizable because the resulting time-domain function, sample function, is not causal. However, by truncating and applying a delay the function could be shifted and, thus satisfy this criteria.
	
	Ideal filter function $H_{d}(\omega)$ has a instantaneous transition at $\omega_c$  resulting in a infinite time response with a decaying factor or a tail that never fully decays to zero. When a real filter is implemented, $h_{d}(n)$ is truncated resulting in some values of the decaying tail left off. This introduces a ringing close to $\omega_c$ known as Gibbs phenomenon. Note that no matter how many taps M are taken, the overshoot stays at a fixed 9\% value of the jump height. Since truncation is therefore unavoidable, the real filter function must have a finite transition (transition-band) between the desired frequencies (pass-band) and the attenuated frequencies (stop-band).
	
	%%%%%%%%%%%%
	%if (ontime) {make pretty diagram}
	%%%%%%%%%%%%

	 Filters can be also classified in two groups according to the response they produce in time-domain: Finite impulse response (FIR) and Infinite impulse response (IIR). 
	
	An IIR filter is defined mathematically as a discrete time function  
	\[y(n) = \sum_{k=0}^{M} b_n x(n-k) + \sum_{k=1}^{N} a_n y(n-k)\] 
	
	By nature the filter is recursive because of the term  $\sum_{k=1}^{N} a_n y(n-k)$ implying the existence of a feedback path that could lead to unstable behavior due the presence of poles in Z-transform transfer function. However, they are widely used because of the reduction in needed coefficients to obtain a narrower transition-band that represent less memory usage and faster responses. 
	  
	An FIR filter is represented as 
	\begin{equation}
		y(n) = \sum_{k=0}^{M-1} b_k x(n-k) 
	\end{equation}
	
	Notice that there is no recursive term which lead to a always stable function or alternatively, no poles in Z-transform transfer function. This means that the filter only depends on current and past inputs. 
	
	Saying this, the analysis of the chosen topology, FIR Low-pass (LP) filter begins by looking the response of feeding the FIR equation with a impulse of $x(n) = \delta(n)$ results in  
	
	\[ y(n) = \sum_{k=0}^{M-1} b_k \delta(n-k) = b_n\]
	
	 $\delta(n-k)$ is zero except when n=k, the impulse response are the coefficients itself. This is the unit sample response $h(n)$, so, by notation (1) becomes
	 \begin{equation}
	 	y(n) = \sum_{k=0}^{M-1} h(k) x(n-k) 
	 \end{equation}
	
	So the problem is reduced to determine the coefficients $h(n)$ for a desired frequency response $H_d(\omega)$. Before moving, another condition must be satisfied to provide correct characterization for a LP filter: if the symmetry condition 
	\[h(n) = h(M-1-n) \]
	
	 is satisfied the filter has a linear-phase or a constant group delay, ensuring the filter doesn't distort the signal.
	 
	 Still one detail remain, and its that the resultant $h_d(n)$, from applying inverse DTFT to $H_d(\omega)$, is infinite in duration and must be truncated to a $n=M-1$, that is archive though Window Method, resulting in
	 \[h(n) = h_d(n) w(n) \] 
	 
	 h(n) is a multiplication in time. In frequency this domain yields the convolution of $H_d(\omega)$ and $W(\omega)$, which function's spectrum has a main lobe that produces a transition band. Also this spectrum has sides lobes that produces ripples in pass and stop band. 
	 %Kinda end of the introduction to the introduction 
	
	The following $H_d(\omega)$ is proposed*, 
	\begin{equation}
		H_d(\omega)=
		\begin{cases}
			e^{-j\omega\frac{M-1}{2}}, & 0 \leq |\omega| \leq \omega_c,\\[6pt]
			0, & \text{otherwise}.
		\end{cases}
	\end{equation}
	 
	 which magnitude response follows the ideal rectangular window( keeping the amplitude of the signal in the pass-band) and, which frequency response presents a $\frac{M-1}{2}$ delay making it a causal function.
	 
	 Applying inverse DTFT  to (3) 
	 
	 \begin{equation*}
	 	\begin{aligned}
	 		h_d(n)
	 		&=
	 		\frac{1}{2\pi}
	 		\int_{-\omega_c}^{\omega_c}
	 		e^{j\omega\left(n-\frac{M-1}{2}\right)}
	 		\,d\omega =
	 		\frac{
	 			\sin\left[
	 			\omega_c\left(n-\frac{M-1}{2}\right)
	 			\right]
	 		}{
	 			\pi\left(n-\frac{M-1}{2}\right)
	 		},
	 		\qquad
	 		n\neq\frac{M-1}{2}.
	 	\end{aligned}
	 \end{equation*}
	
	where $\omega_c$ is the normalized cutoff frequency that satisfies $0<\omega_c < \pi$. 
	
	Notice that, for even M numbers the function never indeterminate because n are integer values. 
	
	By selecting a Hamming window, which has a wide main-lobe and low side-lobes,
	
	\begin{equation*}
		w(n)=0.54-0.46\cos\left(\frac{2\pi n}{M-1}\right),
		\qquad 0\leq n\leq M-1.
	\end{equation*}
	
	the ideal magnitude response of $H_d(\omega)$ acquire a wide transition-band and small ripples in pass and stop bands.
	
	With a even number of tap and the symmetric condition, at Nyquist frequency $\omega = \pi$, 
	
	\[	H(\omega) = \sum_{n=0}^{M-1} h(n) e^{-j\omega n} \]
	\[	H(\pi) = \sum_{n=0}^{M-1} h(n) e^{-jn\pi} =  \sum_{n=0}^{M-1} h(n) (-1)^n \]
	
	Taking as a example M=4, 
	\[	H(\pi) = h(0)-h(1)+ h(2)-h(3)\]
	Applying the symmetry condition yields $h(0)=h(3) $ and $h(1)-h(2)$, so 
	\[H(\pi) = 0\]
	
	Its easy to see that coefficients cancels for all M evens because the interspersed signs and the symmetry of them, assuring that at the highest useful frequency, the filter suppresses the signal. 

	With both expressions established, to avoid tedious manual calculations, the computation of the filter coefficients was performed in Python.
	
	
		
	% \alert{rewrite this paragraph}
\end{document}